\documentclass[RNAAS]{aastex631}

\begin{document}

\title{Candidate Rotation Periods for 34 Main-belt Asteroids from Gaia DR3 Sparse Photometry}

\author{Oleg Tropynin}
\affiliation{Independent researcher, Kazakhstan}

\begin{abstract}
I searched Gaia DR3 sparse photometry of 598 well-observed numbered asteroids for rotation periods,
fitting at each trial frequency a linear model with season offsets, a linear phase-angle term and a
two-harmonic Fourier series. For confident results whose period was not used to tune the method, 96 of
122 (79\%) agree within 3\% with periods published by others; accuracy is low above 50\,h, where the
method usually returns half of the true period. I list 34 asteroids whose confident period was absent
from all major period compilations, as targets for confirmation. A second analysis of the same data with
a refined method finds the same period for all 34.
\end{abstract}

\section{Data}

I ranked the numbered asteroids in the Gaia DR3 \texttt{sso\_source} table \citep{Tanga2023} by their
number of observations and kept the 6{,}000 best observed. Going down this list, I took the first 300
that have a rotation period of 2--100\,h in JPL SBDB, which follows the October 2023 LCDB
\citep{Warner2009} (\emph{validation set}), and the first 298 that have no JPL period (\emph{target set}).

Gaia reports one $G$ magnitude per CCD crossing in \texttt{sso\_observation}. For each transit
(\texttt{transit\_id}) I took the unweighted means of \texttt{epoch} (TCB) and \texttt{g\_mag} over its CCD
observations; $G_i$ is this mean. Its uncertainty is
$\sigma_i=\sqrt{(1.0857\,\max\sigma_F/F)^2+0.03^2}$, with the maximum taken over the CCDs. JPL Horizons
(observer code 500@-139479, Gaia), queried at the transit epochs, gave the heliocentric distance $r_i$ and
the distance from Gaia $\Delta_i$ (au) and the phase angle $\alpha_i$ (deg). Epochs $t_i$ (days) were
corrected for light time, $t_i\to t_i-\Delta_i/c$. A new \emph{season} starts wherever consecutive transits
are more than 60\,d apart; no minimum number of transits per season is imposed.

\section{Method}

For a trial frequency $f$ (cycles per day) the reduced magnitude is modelled as
\begin{equation}
G_i-5\log_{10}(r_i\Delta_i)=c_{s(i)}+\beta\,\alpha_i+\sum_{k=1}^{2}\left[a_k\cos(2\pi k f t_i)+b_k\sin(2\pi k f t_i)\right],
\end{equation}
where $s(i)$ is the season of transit $i$. All coefficients are free and linear: one offset $c_s$ per
season, which absorbs changes of aspect between seasons, a common phase slope $\beta$ (mag\,deg$^{-1}$),
and four Fourier amplitudes. They are solved by least squares with weights $\sigma_i^{-2}$, giving
$\chi^2(f)$. No correction for the motion of the phase-angle bisector is applied, so periods are
synodic. Before the search, the model without Fourier terms is fitted, transits whose residual exceeds
$5\times1.4826$ times the median absolute deviation of the residuals are removed, and fit and cut are
repeated once.

$\chi^2(f)$ is evaluated on a uniform grid from 0.24 to 12 cycles per day with a step of $0.1/T$, where $T$
is the time span. The grid minimum $f_1$ gives $P_1=1/f_1$. Because most lightcurves are double-peaked and
can be fitted by the second harmonic at half the period, I adopt $P=2P_1$ when $\chi^2(f_1/2)$, computed
exactly, is below $1.5\,\chi^2(f_1)$ and $2P_1\le100$\,h; otherwise $P=P_1$. The factor 1.5 maximises the
number of validation periods within 3\% of the reference ($|P-P_{\rm ref}|/P_{\rm ref}<0.03$; half or double
counts as disagreement) over factors 1.00--3.00 in steps of 0.05. The amplitude is the peak-to-peak range of
the Fourier terms fitted at $1/P$.

The confidence score is $S=\ln(\chi^2_0/\chi^2_1)\,(\chi^2_2/\chi^2_1-1)$, with $\chi^2_0$ the fit without
Fourier terms, $\chi^2_1=\chi^2(f_1)$ whether or not the period is doubled, and $\chi^2_2$ the lowest grid
value outside $\pm5/T$ of $f_1$, $f_1/2$ and $2f_1$ (of $f_1/4$ instead of $2f_1$ when
$\chi^2(f_1/2)<1.25\,\chi^2(f_1)$). Sorting the validation results by decreasing $S$, the threshold
$S\ge0.1604$ is the score of the last asteroid at which the cumulative agreement fraction is still at
least 80\% (from the tenth asteroid on). Code and the list of VizieR tables are archived with the data
\citep{Tropynin2026}.

\section{Validation}

With the doubling factor chosen on one half of the validation set (alternate asteroids in the ranked list)
and tested on the other, 203 of 300 periods agree. Above the threshold:
\begin{itemize}
\item validation set: 172 of 214 agree (in-sample, because this set fixed the threshold);
\item target-set asteroids solved from the same Gaia data with a different method by \citet{Durech2023}:
81 of 97 agree;
\item the other confident target-set asteroids with a rotation period in \citet{Cellino2024},
\citet{MacLennan2026} or \citet{McNeill2023}: 15 of 25 agree.
\end{itemize}
Out of sample, 96 of 122 (79\%) agree. By period, 165 of 203 (2--10\,h), 79 of 89 (10--30\,h), 15 of 17
(30--50\,h) and 10 of 28 (50--100\,h) agree, counting all three groups.

\section{The 34 candidates}

Of the 298 target-set asteroids, 165 pass the threshold. I removed the 97 solved by \citet{Durech2023}, the
25 above, and 35277, whose only published period, 23.61\,h \citep{Liberato2024}, comes from astrometric
wobble but is close to ours (23.88\,h). None of the remaining 42 has a rotation period, in October 2026, in
JPL SBDB, LCDB~V4.0, DAMIT \citep{Durech2010}, the SsODNet ssoBFT table \citep{Berthier2023}, 179 VizieR
tables (5.1 million rows), or the Minor Planet Bulletin volumes 50--53, except 11857, for which
\citet{Durech2022} give 19.72\,h from ATLAS photometry. A second analysis of all of Gaia DR3
\citep{Tropynin2026b} adds a correction for the phase-angle bisector, a calibrated probability that a period
is correct, and searches to 1000\,h. I dropped the seven candidates where its best period differs by more
than 3\%: four where it is twice ours and three where it is a slow solution. The remaining 34
(Table~\ref{tab:1}) are main-belt asteroids with 58--91 transits in 4--7 seasons; the best period of the
second analysis is within 0.2\% of ours for all of them, with probability above 0.8 for 18. Formal period
uncertainties are not computed. Out of sample, about four in five periods below 50\,h agree with published
values. I offer the list as targets for dense-lightcurve confirmation.

\begin{deluxetable*}{rrcrc|rrcrc}
\tablecaption{The 34 candidate rotation periods.\label{tab:1}}
\tablehead{\colhead{Number} & \colhead{$P$ (h)} & \colhead{Amp.} & \colhead{$N$} & \colhead{Q} &
\colhead{Number} & \colhead{$P$ (h)} & \colhead{Amp.} & \colhead{$N$} & \colhead{Q}}
\startdata
5897 & 6.0729 & 0.39 & 67 & A & 44204 & 4.7663 & 0.72 & 65 & A \\
8751 & 3.2402 & 0.25 & 70 & A & 45726 & 7.9396 & 0.33 & 69 & B \\
11621 & 33.7373 & 0.60 & 59 & A & 47917 & 4.7480 & 0.38 & 80 & A \\
11818 & 2.7435 & 0.17 & 66 & B & 50588 & 56.9297$^a$ & 0.55 & 65 & A \\
13363 & 8.9320 & 0.26 & 68 & A & 52653 & 3.7362 & 0.14 & 65 & -- \\
16916 & 3.8301 & 0.15 & 69 & B & 55335 & 5.0076 & 0.14 & 61 & -- \\
17337 & 4.9700 & 0.21 & 65 & B & 56994 & 14.7458 & 0.39 & 64 & A \\
22427 & 5.4857 & 0.46 & 66 & A & 63007 & 4.4476 & 0.33 & 65 & A \\
24359 & 16.1525 & 0.22 & 62 & A & 64303 & 8.8781 & 0.19 & 63 & B \\
24711 & 6.4557 & 0.75 & 61 & B & 65213 & 3.8679 & 0.34 & 73 & A \\
30068 & 8.3416 & 0.27 & 72 & B & 72885 & 3.0007 & 0.32 & 62 & B \\
31611 & 5.4767 & 0.38 & 62 & A & 84508 & 3.0394 & 0.23 & 69 & A \\
32778 & 66.9411$^a$ & 0.56 & 63 & A & 87249 & 4.4550 & 0.19 & 70 & B \\
35653 & 30.0890 & 0.13 & 66 & -- & 88237 & 52.2814$^a$ & 0.10 & 58 & -- \\
40172 & 5.1295 & 0.54 & 74 & A & 91539 & 5.6743 & 0.42 & 91 & A \\
42087 & 81.9684$^a$ & 0.22 & 66 & B & 97578 & 2.8696 & 0.41 & 86 & A \\
44017 & 3.0915 & 0.20 & 65 & B & 105992 & 8.0425 & 0.19 & 58 & B \\
\enddata
\tablecomments{$P$: adopted period after the doubling rule. Amp.: peak-to-peak amplitude of the Fourier terms (mag).
$N$: transits used after outlier removal. Q: tier in the second analysis \citep{Tropynin2026b}, A for
probability $\ge0.8$, B for 0.38--0.8, -- below. $^a$Above 50\,h; the true period may be twice the value given.}
\end{deluxetable*}

\begin{acknowledgments}
This work used data from the ESA mission Gaia, processed by the Gaia DPAC; JPL Horizons and the JPL
Small-Body Database; and VizieR (CDS, Strasbourg). The analysis was designed, coded and checked with the
help of an AI assistant (Claude, Anthropic) under the author's direction.
\end{acknowledgments}

\begin{thebibliography}{}
\bibitem[Berthier et al.(2023)]{Berthier2023} Berthier, J., Carry, B., Mahlke, M., \& Normand, J. 2023, A\&A, 671, A151
\bibitem[Cellino et al.(2024)]{Cellino2024} Cellino, A., et al. 2024, A\&A, 687, A277
\bibitem[\v{D}urech \& Hanu\v{s}(2023)]{Durech2023} \v{D}urech, J., \& Hanu\v{s}, J. 2023, A\&A, 675, A24
\bibitem[\v{D}urech et al.(2022)]{Durech2022} \v{D}urech, J., V\'avra, M., Van\v{c}o, R., \& Erasmus, N. 2022, FrASS, 9, 809771
\bibitem[\v{D}urech et al.(2010)]{Durech2010} \v{D}urech, J., Sidorin, V., \& Kaasalainen, M. 2010, A\&A, 513, A46
\bibitem[MacLennan et al.(2026)]{MacLennan2026} MacLennan, E. M., et al. 2026, A\&A, 707, A131
\bibitem[Liberato et al.(2024)]{Liberato2024} Liberato, L., Tanga, P., Mary, D., et al. 2024, A\&A, 688, A50
\bibitem[McNeill et al.(2023)]{McNeill2023} McNeill, A., et al. 2023, AJ, 166, 152
\bibitem[Gaia Collaboration et al.(2023)]{Tanga2023} Gaia Collaboration, Tanga, P., et al. 2023, A\&A, 674, A12
\bibitem[Tropynin(2026a)]{Tropynin2026} Tropynin, O. 2026a, Candidate rotation periods for main-belt asteroids from Gaia DR3 sparse photometry, Zenodo, v1.3.0, doi:\dataset[10.5281/zenodo.23624115]{https://doi.org/10.5281/zenodo.23624115}
\bibitem[Tropynin(2026b)]{Tropynin2026b} Tropynin, O. 2026b, Rotation periods and spin senses of asteroids from Gaia DR3 sparse photometry, Zenodo, doi:\dataset[10.5281/zenodo.23638281]{https://doi.org/10.5281/zenodo.23638281}
\bibitem[Warner et al.(2009)]{Warner2009} Warner, B.~D., Harris, A.~W., \& Pravec, P. 2009, Icarus, 202, 134
\end{thebibliography}

\end{document}
